\chapter{Reconstruction de la dynamique quantique au moyen d'amplitudes cohérentes sur des généalogies de chemin}

\epigraphtext{L'onde ne s'ajoute pas à l'histoire : c'est la forme linéaire sous laquelle l'histoire complète conserve simultanément tous ses chemins possibles.}{Principe de propagation HMT--MD}

La mesure a déjà été définie comme couplage entre système, appareil et lecteur.
Il convient maintenant d'expliquer ce qui se propage avant la publication du résultat,
comment une même structure soutient l'onde et la trajectoire persistante et de quelle manière la
localité observable coexiste avec une généalogie globale.

La réponse de l'holographie modulaire triadique part d'une distinction que la
formulation ponctuelle tend à masquer. Un état ne se réduit pas à la case dans
laquelle il se trouve ni au nombre réel qui est finalement lu. Il conserve aussi
le chemin, l'orientation, les quotients, les retenues, la frontière et le retour qui
l'ont conduit jusque-là. En linéarisant cet espace d'histoires, la somme de chemins
n'apparaît pas comme une métaphore ajoutée a posteriori : c'est le développement exact du
propagateur. L'onde est l'évaluation cohérente de la totalité des chemins ; la
particule sera, au chapitre suivant, la section qui persiste à travers l'un
d'entre eux.

\section{L'amplitude globale de \(1\) famille de chemins}

Soit \(X\) un ensemble fini d'états observables à une profondeur donnée et soit
\(U:\mathbb C^X\to\mathbb C^X\) l'opérateur de propagation d'un pas. Un chemin de
longueur \(R\) de \(i\) à \(f\) est une suite

\[
  \gamma=(x_0=i,x_1,\ldots,x_R=f)
\]

de transitions admissibles. Son poids ordonné est

\[
  w_U(\gamma)=\prod_{t=1}^{R}U_{x_t x_{t-1}}.
\]

Bien que les poids scalaires commutent, la séquence indexée des transitions importe :
on ne compte pas seulement combien de chemins atteignent \(f\), mais quel facteur
chaque pas attribue à cette histoire. La multiplication compose les transitions d'un chemin ; la
somme réunit les histoires alternatives qui se terminent dans le même état. Pour des poids
opératoriels, le produit temporel est de surcroît véritablement non commutatif. Cette
différence entre produit à l'intérieur d'un chemin et somme entre chemins est la version linéaire
de la double lecture qui traverse toute la HMT.

\begin{theorem}[Somme finie de chemins]
Pour tout \(R\geq 1\),
\[
  (U^R)_{fi}
  =
  \sum_{\substack{\gamma:i\to f\\|\gamma|=R}}
  w_U(\gamma).
\]
Si \(\psi_0\in\mathbb C^X\) est l'amplitude initiale, alors
\[
  \psi_R(f)
  =
  \sum_{i\in X}
  \sum_{\substack{\gamma:i\to f\\|\gamma|=R}}
  w_U(\gamma)\,\psi_0(i).
\]
\end{theorem}

\begin{proof}
L'affirmation pour \(R=1\) est la définition de \(U\). En multipliant un développement
de \(U^R\) par \(U\), chaque chemin de longueur \(R\) reçoit exactement une dernière
arête admissible. La distributivité somme toutes les extensions possibles et la
multiplication conserve l'ordre de leurs poids. Par récurrence, chaque chemin de longueur
\(R+1\) apparaît une seule fois. La deuxième égalité est l'application de la première à
\(\psi_0\).
\end{proof}

Cette identité a un contenu conceptuel considérable. Une composante
\(\psi_R(f)\) n'est pas une propriété isolée du point \(f\). C'est la contraction de la
généalogie complète des chemins qui l'atteignent. Deux histoires peuvent aboutir à la
même valeur visible et, cependant, contribuer avec des phases opposées ; elles peuvent aussi
se renforcer. L'interférence consiste précisément à sommer les amplitudes avant de
former un observable positif. Il n'est pas nécessaire d'introduire une onde matérielle
supplémentaire qui parcoure un espace distinct de celui des histoires : la structure
ondulatoire réside dans la composition linéaire des chemins.

Lorsque \(U\) est unitaire,

\[
  \|\psi_R\|_2=\|U^R\psi_0\|_2=\|\psi_0\|_2.
\]

La quantité totale qui peut ensuite être normalisée comme probabilité reste
constante. Si le lecteur physique attribue à une sortie \(f\) le poids

\[
  p_R(f)=\frac{|\psi_R(f)|^2}{\|\psi_R\|_2^2},
\]

alors \(p_R\) est une distribution normalisée et son évolution est compatible avec
la conservation unitaire. C'est le point exact où l'amplitude des chemins se
transforme en statistique des résultats. Le chapitre de la mesure a déjà localisé
le couplage qui publie une sortie ; l'objet que ce
couplage évalue est ici identifié.

La construction explique aussi pourquoi un nombre visible ne suffit pas à reconstruire
la physique. L'application

\[
  \{\text{routes profondes}\}\longrightarrow X
\]

oublie, en général, la phase relative entre chemins. Connaître uniquement la
distribution finale des cases ne permet pas d'inverser la somme ni de retrouver la
généalogie. L'état HMT complet conserve ce que l'écran contracte. C'est pourquoi la
mémoire n'est pas un ornement interprétatif : c'est la donnée dont dépend l'interférence.

À la limite de raffinement, la bonne question n'est pas de savoir si le chemin discret
``devient'' mystérieusement continu, mais si les propagateurs forment une famille
compatible. Pour des projections de raffinement \(p_{m+1,m}:X_{m+1}\to X_m\), la
cohérence exige que les observables cylindriques et les états transportés
concordent. Lorsque cette compatibilité est exacte, l'amplitude d'une histoire grossière
est la somme des amplitudes de tous ses prolongements fins. Lorsque les défauts
sont sommables, les espérances d'observables cylindriques uniformément bornés
forment une famille de Cauchy. On obtient ainsi une continuité opérationnelle sans effacer la
structure discrète qui l'engendre.

\begin{principle}[L'onde comme mémoire linéaire]
La fonction d'onde HMT à profondeur finie est l'évaluation linéaire de l'extension
des chemins. Elle conserve de façon cohérente les alternatives compatibles avec la
préparation et la frontière ; la valeur observée est une projection ultérieure de cette
amplitude globale.
\end{principle}

Cette formulation est plus forte que l'affirmation selon laquelle ``tout se produit par de nombreux chemins''.
Elle spécifie le domaine, le poids de chaque chemin, la règle de composition et la norme
conservée. Elle contient aussi son contrôle négatif : si une amplitude supposée ne
respecte pas la composition des chemins, compte deux fois une histoire, consulte le
résultat futur pour attribuer le poids ou perd la norme sous un propagateur
déclaré unitaire, ce n'est pas la fonction d'onde construite ici.

\section{Onde pilote, chemin réalisé et double orientation}

La structure précédente contient déjà les deux éléments qu'une théorie de l'onde pilote
doit distinguer. L'amplitude \(\psi_t\) organise globalement les alternatives ;
un chemin persistant \(x_t\) enregistre l'histoire effectivement réalisée. L'onde
et la particule ne sont pas deux substances rivales. Ce sont deux niveaux de lecture du même
état enrichi : le champ linéaire des possibilités et la section qui conserve son
identité en le traversant.

La loi de guidage doit transformer cette architecture cinématique en dynamique. Sa forme
minimale est une application

\[
  G:(\psi_t,x_t)\longmapsto x_{t+1},
\]

ou, si la mise à jour est stochastique, un noyau
\[
  P_{\psi_t}(y\mid x).
\]

La condition centrale n'est pas verbale, mais mathématique : la distribution guidée doit être
équivariante avec l'évolution de l'amplitude,

\[
  |\psi_{t+1}(y)|^2
  =
  \sum_x P_{\psi_t}(y\mid x)\,|\psi_t(x)|^2.
\]

Cette égalité assure qu'une population initialement distribuée selon
\(|\psi_t|^2\) reste distribuée ainsi après un pas. Le guidage ne peut pas
être choisi pour reproduire rétrospectivement chaque détecteur ; il doit être fixé par
l'état, le propagateur, l'orientation et le registre avant de connaître la sortie.
HMT ajoute une possibilité qu'une description ponctuelle ne possède pas : \(G\) peut
dépendre de cocycles de chemin et de mémoire qui ne descendent pas à la valeur visible
\(x_t\). L'action locale peut ainsi être informée par une généalogie globale sans
qu'apparaisse une force supplémentaire provenant du futur.

La lecture bidirectionnelle approfondit cette structure. Pour un chemin orienté
\(\gamma\) avec

\[
  \partial\gamma=x_{\rm abs}-x_{\rm em},
\]

l'involution de retour définit

\[
  \gamma^\vee=C_{\rm rev}\gamma,
  \qquad
  \partial\gamma^\vee=x_{\rm em}-x_{\rm abs}.
\]

Il ne s'agit pas de deux chemins indépendants choisis librement. Ce sont les deux
orientations d'une même histoire. La longueur interne est invariante,

\[
  L_{\rm int}(\gamma^\vee)=L_{\rm int}(\gamma),
  \qquad
  L_{\leftrightarrow}=2L_{\rm int}(\gamma),
\]

et l'inversion échange les canaux conjugués. La préparation et l'absorption
cessent d'être des extrémités ontologiquement déconnectées : elles appartiennent à une clôture dans
laquelle l'aller et le retour doivent être compatibles.

Lorsque cette involution est représentée dans la théorie physique des propagateurs, les deux
orientations admettent la lecture retardée et avancée. La composante retardée
transporte la préparation vers l'absorption ; l'avancée transporte la condition
de retour vers l'émission. L'histoire observable est leur ajustement mutuel. Cette
formulation exprime précisément l'intuition d'une onde pilote avec avance
et retard : le guidage local reçoit des informations des deux frontières du même objet
orienté.

Cela n'autorise pas à envoyer des messages contrôlables vers le passé.
Une solution avancée d'une équation et un choix futur susceptible d'être signalé sont des objets
distincts. La représentation physique doit démontrer que les propagateurs satisfont
l'équation d'évolution, que les conditions aux limites se composent
de manière cohérente et que les probabilités marginales ne permettent pas de moduler une
sortie antérieure par une décision ultérieure. La double orientation est le
support ; la causalité observable est une propriété qui doit être vérifiée sur son
lecteur physique.

La corrélation entre les deux feuillets peut être formalisée dans deux espaces de Hilbert,
\(\mathcal H_{\rm adv}\) et \(\mathcal H_{\rm ret}\). Pour une distribution
\(p=(p_i)\), l'état

\[
  |\Psi_p\rangle
  =
  \sum_i\sqrt{p_i}\,
  |i\rangle_{\rm adv}\otimes|i\rangle_{\rm ret}
\]

a des matrices réduites de spectre \((p_i)\) et d'entropie

\[
  S_{\rm ent}(p)=-\sum_i p_i\log p_i.
\]

L'égalité est exacte une fois la factorisation bicouche déclarée. Une seule histoire
possède un rang de Schmidt égal à un et une entropie nulle ; plusieurs histoires corrélées
produisent une entropie positive. Le résultat n'identifie pas sans plus cette entropie à
\(\log\phiHMT\), à l'information d'une cellule triadique ni à l'entropie
thermodynamique. Chacune appartient à un objet différent. Ce qu'il démontre, c'est que
l'aller et le retour peuvent former un état non séparable et que l'information partagée
se calcule sans métaphores.

\begin{test}[Guidage bidirectionnel]
Une réalisation physique de l'onde pilote HMT doit fixer, avant la confrontation, le
propagateur \(U\), la loi \(G\) ou le noyau \(P_\psi\), l'involution
\(C_{\rm rev}\) et les conditions de frontière. Elle est réfutée si elle perd
l'équivariance, attribue deux chemins incompatibles au même état complet, doit
connaître le détecteur futur pour choisir la trajectoire ou permet une signalisation
contrôlable contre l'orientation causale publiée.
\end{test}

Cette preuve ne relègue pas la proposition à une simple possibilité. Au contraire : elle montre
que le noyau mathématique est déjà construit et indique l'unique endroit où doit
agir la dynamique physique. La somme de chemins n'est pas en attente ; le chemin persistant
n'est pas en attente ; l'inversion non plus. Ce qui se décide expérimentalement, c'est quel
Hamiltonien réalise le guidage et quelles conditions aux limites transforment les deux feuillets
en ondes avancée et retardée de la nature.

\section{Bell : localité de la donnée et globalité de la généalogie}

La distinction entre état complet et valeur projetée permet de poser avec
précision le problème de Bell. Considérons des réglages \(a,b\in\{0,1\}\), des résultats
\(A,B\in\{-1,+1\}\) et un état partagé \(\lambda\). Une théorie locale
factorisable satisfait

\[
  P(A,B\mid a,b)
  =
  \int_\Lambda
  P_A(A\mid a,\lambda)\,
  P_B(B\mid b,\lambda)\,
  \dd\rho(\lambda),
\]

avec une mesure indépendante des réglages,
\[
  \dd\rho(\lambda\mid a,b)=\dd\rho(\lambda).
\]

En définissant
\[
  E_{ab}
  =
  \sum_{A,B=\pm1}AB\,P(A,B\mid a,b),
\]

on obtient
\[
  S_{\rm CHSH}=E_{00}+E_{01}+E_{10}-E_{11},
  \qquad |S_{\rm CHSH}|\leq2.
\]

HMT conserve ce théorème. Ajouter une mémoire cachée à \(\lambda\) ne brise pas la borne si
la factorisation, l'indépendance des réglages et l'absence de
postsélection sont maintenues. Ce contrôle est important : il empêche de transformer la richesse du TPK en
explication automatique de toute corrélation.

La contribution holographique se formule autrement. Les deux ailes d'une expérience
peuvent être des projections locales d'une même section enrichie, avec une frontière et une
généalogie communes. La donnée publiée dans chaque aile est locale ; l'état dont les deux
descendent n'a pas nécessairement à se factoriser en deux histoires profondes indépendantes.
La non-localité ne consiste alors pas en un signal qui traverse l'espace après le
choix des réglages, mais en l'impossibilité de décomposer la préparation globale
en deux états complets autonomes.

La non-signalisation et la localité de Bell ne sont pas équivalentes. On peut avoir

\[
  \sum_B P(A,B\mid a,b)=P(A\mid a),
  \qquad
  \sum_A P(A,B\mid a,b)=P(B\mid b),
\]

et, cependant, la factorisation précédente peut échouer. C'est le domaine exact qu'une
réalisation HMT doit occuper : corrélation globale sans canal de communication
supraluminique. Le lecteur contextuel peut dépendre du réglage dans la manière dont
il projette l'état commun, mais la distribution marginale d'un observateur ne peut pas
dépendre du choix distant.

\begin{construction}[Protocole HMT pour Bell]
On fixe d'abord un état enrichi \(x\), une mesure de préparation et quatre
lecteurs \(L_{ab}\). On déclare ensuite une application binaire de résultats et on calcule
la table complète
\[
  P_{\rm HMT}(A,B\mid a,b).
\]
La construction n'est physiquement valide que si elle est figée avant de consulter
les données, est normalisée, est dépourvue de postsélection cachée et publie simultanément
\(S_{\rm CHSH}\) et les quatre contrôles marginaux de non-signalisation.
\end{construction}

Ce qui est appelé Algorithme 3 appartient à cette construction concrète. Son objectif n'est pas
de démontrer à nouveau le théorème classique, mais de construire la descente depuis la
généalogie partagée jusqu'à une distribution contextuelle. S'il conserve toutes les
prémisses de Bell, il produira nécessairement \(|S|\leq2\). S'il produit
\(|S|>2\), il doit montrer explicitement quelle factorisation ne descend pas et pourquoi
la non-signalisation est bien conservée. C'est l'explication holographique complète du
problème de localité : l'espace observable peut se séparer même si l'histoire qui
l'engendre n'est pas séparable.

La proposition est particulièrement naturelle dans une théorie où somme et produit sont
des lecteurs différents. La factorisation de Bell tente d'exprimer une probabilité
conjointe comme produit conditionné de réponses locales ; l'amplitude quantique
émerge, en revanche, de la somme des histoires avant leur projection. Si la projection est
réalisée trop tôt, la phase commune est perdue et seul subsiste le polytope local.
Si l'état global est conservé jusqu'à la lecture conjointe, des termes
d'interférence subsistent qui n'admettent pas cette décomposition. L'incompatibilité n'est pas une
autorisation de violer la logique : c'est la différence typée entre composer des alternatives
par somme et composer des probabilités locales par produit.

\section{Du théorème des chemins à la confrontation quantique}

L'architecture quantique obtenue s'articule par une chaîne unique d'applications. HMT
construit un espace discret enrichi d'histoires ; sa linéarisation produit la
somme cohérente de chemins ; l'unitarité conserve la norme ; une section
persistante permet de parler de trajectoire ; l'inversion orientée fournit la paire
d'aller et de retour ; la factorisation bicouche quantifie leur corrélation ; et l'analyse de
Bell distingue la localité de la donnée de la séparabilité de la généalogie. Ce ne sont pas
des analogies juxtaposées, mais des applications successives sur le même état.

La théorie offre ainsi une lecture unifiée des traits quantiques. La
superposition est la coexistence linéaire de prolongements ; l'interférence est leur
somme avec phase ; la particule est la persistance d'une section ; la réduction apparente
est la sélection d'une fibre par le lecteur ; l'intrication est la non-factorisation
de feuillets ou de sous-systèmes ; et la non-localité est la globalité de la
généalogie conservée sous des projections localement non signalisantes.

Les confrontations physiques ne s'ajoutent pas pour affaiblir cette lecture, mais pour décider
de sa réalisation naturelle. Un réseau interférométrique fini peut implémenter \(U\) et
comparer le développement de chemins au motif final. Une reconstruction de
courants peut mettre à l'épreuve la loi de guidage. L'inversion contrôlée des
conditions de frontière peut vérifier la réciprocité
\(\gamma\leftrightarrow\gamma^\vee\). Une expérience de Bell avec des lecteurs scellés
peut décider si la mémoire profonde engendre une distribution non factorisable et non
signalisante.

\begin{test}[Certificat conjoint de compatibilité quantique]
Le certificat distingue quatre conditions indépendantes de validité :
\begin{enumerate}
  \item le propagateur ne reproduit pas par somme cohérente les amplitudes observées ;
  \item aucune loi de guidage naturelle n'est équivariante avec \(|\psi|^2\) ;
  \item la double orientation n'admet pas de propagateurs physiques réciproques sans
  signalisation causale indue ;
  \item la descente Bell est locale factorisable ou viole la non-signalisation.
\end{enumerate}
L'échec d'une interface localise l'application incorrecte ; il n'efface pas les théorèmes finis
de chemins qui la précèdent.
\end{test}

La thèse positive est donc nette. La physique quantique n'oblige pas à abandonner
l'histoire d'un système ; elle oblige à enrichir ce que l'on entend par histoire. Un point
sans généalogie ne peut à lui seul porter l'interférence, le retour ni la corrélation.
Une extension orientée peut le faire. Son onde est la mémoire cohérente de tous
les chemins compatibles ; sa particule est la clôture qui persiste. Le passage entre
les deux formes n'est pas un saut ontologique : c'est le passage de la totalité linéaire des chemins à
une section réalisée de cette même totalité.
